% Source PDF: source/普通高考/2013/2013江西文.pdf, page 1, question 7.
% Original PNG reference: img/2013/jiangxi_liberal/q7_fig1.png.
% Node -> directed-edge list: 开始→(i=1,S=0)→i=i+1→i是奇数;
% 否→S=2*i+1→第二判定；是→S=2*i+2→第二判定；
% 第二判定是→回到 i=i+1 前的主干；否→输出 i→结束。
% The two assignment branches enter the same merge with independent arrows;
% the loop return has its own arrow into the pre-increment merge.
\begin{tikzpicture}[
  x=1cm,y=1cm,
  >=Stealth,
  line width=0.55pt,
  line cap=round,line join=round,
  every node/.style={font=\normalsize,anchor=center,align=center,
    execute at begin node={\setlength{\parindent}{0pt}}},
  flowbox/.style={draw,minimum height=.66cm,inner xsep=4pt,inner ysep=2pt,
    anchor=center,align=center},
  io/.style={draw,shape=trapezium,trapezium left angle=78,trapezium right angle=102,
    trapezium stretches=false,minimum height=.76cm,inner xsep=4pt,inner ysep=2pt,
    anchor=center,align=center},
  terminal/.style={flowbox,rounded corners=7pt,minimum width=1.35cm,
    minimum height=.70cm},
  decision/.style={flowbox,diamond,aspect=2.8,minimum width=4.55cm,
    minimum height=1.08cm,inner sep=2pt},
  branchlabel/.style={inner sep=0pt,anchor=center,align=center,
    execute at begin node={\setlength{\parindent}{0pt}}},
]
  \node[terminal] (start) at (0,10.30) {开始};
  \node[flowbox,minimum width=3.35cm] (init) at (0,8.95) {$i=1,S=0$};
  \coordinate (preinc) at (0,8.28);
  \node[flowbox,minimum width=2.55cm] (inc) at (0,7.60) {$i=i+1$};
  \node[decision] (parity) at (0,6.18) {$i\text{ 是奇数}$};
  \node[flowbox,minimum width=2.70cm] (odd) at (4.00,6.18) {$S=2*i+1$};
  \node[flowbox,minimum width=2.85cm] (even) at (0,4.90) {$S=2*i+2$};
  \coordinate (branchmerge) at (0,4.00);
  \node[decision,minimum width=4.60cm] (test) at (0,3.10) {};
  \node[io,minimum width=2.10cm] (output) at (0,1.62) {输出 $i$};
  \node[terminal] (finish) at (0,0.28) {结束};

  % Main chain and the pre-increment merge.
  \draw[->] (start.south) -- (init.north);
  \draw (init.south) -- (preinc);
  \draw[->] (preinc) -- (inc.north);
  \draw[->] (inc.south) -- (parity.north);

  % The parity branches are distinct and merge once before the blank test.
  \draw[->] (parity.east) -- (odd.west);
  \draw[->] (parity.south) -- (even.north);
  % Put each branch arrowhead on its own incoming segment; the merge itself
  % remains a clean, unambiguous junction with no stacked arrowheads.
  \draw[->] (odd.south) -- (4.00,4.30);
  \draw (4.00,4.30) -- (4.00,4.00) -- (branchmerge);
  \draw[->] (even.south) -- (0,4.30);
  \draw (0,4.30) -- (branchmerge);
  \draw[->] (branchmerge) -- (test.north);

  % The second decision loops back to the pre-increment merge or outputs i.
  \coordinate (loop-left) at (-3.55,3.30);
  \coordinate (loop-entry) at (-3.55,8.28);
  \draw[->] (test.west) -- (loop-left) -- (loop-entry) -- (-0.40,8.28);
  \draw (-0.40,8.28) -- (preinc);
  \draw[->] (test.south) -- (output.north);
  \draw[->] (output.south) -- (finish.north);

  \node[branchlabel,anchor=south] at (2.72,6.50) {否};
  \node[branchlabel,anchor=west] at (0.22,5.35) {是};
  \node[branchlabel,anchor=east] at (-2.85,2.88) {是};
  \node[branchlabel,anchor=west] at (0.22,2.52) {否};

  \path[use as bounding box] (-3.90,10.70) rectangle (5.50,0.15);
\end{tikzpicture}
